En un transformador que consta de dues bobines, la bobina primària té $N_\mathrm{p}$ espires i la secundària té $N_\mathrm{s}$ espires.

\begin{apartat}{1}
Deduïu, a partir de la conservació del flux magnètic, la fórmula per a obtenir la tensió del circuit secundari quan connectem la bobina primària d'un transformador a una tensió alterna $\varepsilon$.

Si $N_\mathrm{p} = 1\,200$ espires i $N_\mathrm{s} = 300$ espires, calculeu la tensió eficaç a la bobina secundària quan connectem la bobina primària a una tensió eficaç de 230~V.
\end{apartat}

\begin{apartat}{1}
Calculeu la intensitat eficaç en el circuit primari si pel circuit secundari circulen 2,0~A d'intensitat eficaç. Feu un esquema i indiqueu-hi cada element del transformador, sabent que les dues bobines estan enrotllades sobre un nucli de ferro comú.
\end{apartat}

\nota{Considereu un transformador ideal.}
